% Part: first-order-logic
% Chapter: sequent-calculus
% Section: rules-and-proofs

\documentclass[../../../include/open-logic-section]{subfiles}

\begin{document}

\iftag{FOL}
      {\olfileid{fol}{seq}{rul}}
      {\olfileid{pl}{seq}{rul}}

\olsection{Regels en \usetoken{P}{derivation}}

In het vervolg stellen $\Gamma, \Delta, \Pi, \Lambda$ eindige rijen
!!{sentence}s voor.

\begin{defn}[Sequent]
Een \emph{sequent} is een uitdrukking van de vorm
\[
\Gamma \Sequent \Delta
\]
waarbij $\Gamma$ en $\Delta$ eindige (mogelijk lege) rijen
!!{sentence}s van de taal $\Lang L$ zijn. $\Gamma$ heet het
\emph{antecedent} en $\Delta$ het \emph{succedent}.
\end{defn}

\begin{explain}
Het intuïtieve idee achter een sequent is dat, als alle !!{sentence}s
in het antecedent gelden, ten minste een van de !!{sentence}s in het
succedent geldt. Als $\Gamma = \tuple{!A_1, \dots, !A_m}$ en $\Delta =
\tuple{!B_1, \dots, !B_n}$, dan geldt $\Gamma \Sequent \Delta$ dus dan
en slechts dan als
\[
(!A_1 \land \cdots \land !A_m) \lif (!B_1 \lor \cdots \lor
!B_n)
\]
geldt. Er zijn twee bijzondere gevallen: $\Gamma$ is leeg of $\Delta$
is leeg. Als $\Gamma$ leeg is, dus als $m = 0$, dan geldt $\quad
\Sequent \Delta$ dan en slechts dan als $!B_1 \lor \dots \lor !B_n$
geldt. Als $\Delta$ leeg is, dus als $n = 0$, dan geldt $\Gamma
\Sequent \quad$ dan en slechts dan als $\lnot(!A_1 \land \dots \land
!A_m)$ geldt. Een sequent heet geldig dan en slechts dan als de
bijbehorende !!{sentence} geldig is.
\end{explain}

Als $\Gamma$ een rij !!{sentence}s is, schrijven we $\Gamma, !A$ voor
de rij die ontstaat door $!A$ rechts aan~$\Gamma$ toe te voegen (en
$!A, \Gamma$ voor de rij die ontstaat door $!A$ links aan~$\Gamma$ toe
te voegen). Als ook $\Delta$ een rij !!{sentence}s is, dan is $\Gamma,
\Delta$ de concatenatie van beide rijen.

\begin{defn}[Beginsequent]
Een \emph{beginsequent} is een sequent
\iftag{prvFalse,prvTrue}{van een van de volgende vormen:
  \begin{enumerate}
    \item $!A \Sequent !A$
    \tagitem{prvTrue}{$\quad \Sequent \ltrue$}{}
    \tagitem{prvFalse}{$\lfalse \Sequent \quad$}{}
  \end{enumerate}}
{van de vorm $!A \Sequent !A$} voor elke !!{sentence} $!A$ van de taal.
\end{defn}

!!^{derivation}s in de sequentencalculus zijn bepaalde bomen van
sequenten. De bovenste sequenten zijn beginsequenten. Staat een sequent
onder een of twee andere sequenten, dan moet hij daaruit correct volgen
volgens een afleidingsregel. De regels voor $\Log{LK}$ worden in twee
hoofdsoorten verdeeld: \emph{logische} en \emph{structurele} regels.
Een logische regel wordt genoemd naar de !!{main operator} van de
!!{sentence} met $!A$ en/of $!B$ in de onderste sequent. Elke logische
regel heeft twee versies: een waarmee een sequent wordt afgeleid waarin
de !!{sentence} met de !!{operator} links staat, en een waarin die
!!{sentence} rechts staat.

\end{document}
